\chapter{Boussinesq Yaklaşımı ve Konveksiyon}
\label{ch:boussinesq}

% ============================================================
\section{Boussinesq Yaklaşımı Nedir?}
\label{sec:boussinesq_nedir}

Gerçek akışkanlarda sıcaklık değişimi yoğunluğu değiştirir. Sıcak hava hafifler, soğuk hava ağırlaşır. Tam sıkıştırılabilir denklemler bu etkiyi eksiksiz yakalar, ama çözümü çok pahalıdır (ses dalgalarını da çözmek gerekir).

Joseph Valentin Boussinesq (1903), basit ama güçlü bir yaklaşım önerdi:

\begin{definition}[Boussinesq Yaklaşımı]
Yoğunluk değişimi \textbf{yalnızca kaldırma (buoyancy) teriminde} hesaba katılır. Diğer tüm terimlerde yoğunluk sabit ($\rho_0$) alınır.
\end{definition}

Başka bir deyişle:
\begin{itemize}
    \item Süreklilik denkleminde: $\divergence\vect{u} = 0$ (sanki sıkıştırılamaz).
    \item Momentum denkleminde: Atalet terimi $\rho_0\,D\vect{u}/Dt$, viskoz terim $\mu\laplacian\vect{u}$ --- $\rho = \rho_0$ kullanılır.
    \item \textbf{Tek istisna:} Yerçekimi terimi $\rho\vect{g}$'de gerçek yoğunluk $\rho(T)$ kullanılır.
\end{itemize}

Bu ``tutarsız'' gibi görünebilir, ama fiziksel olarak gerekçelidir: yoğunluk değişimi küçüktür ($\Delta\rho/\rho_0 \ll 1$), atalet ve viskoz terimlerde ihmal edilebilir. Ama yerçekimi kuvveti $\rho g$ büyük bir sayıdır; küçük bir $\Delta\rho$ bile önemli bir kaldırma kuvveti yaratır.


% ============================================================
\section{Boussinesq Denklemlerinin Türetmesi}
\label{sec:boussinesq_turetme}

\textbf{Adım 1: Durum denklemi (lineer yaklaşım).}

Yoğunluğun sıcaklığa bağımlılığını birinci mertebe Taylor açılımıyla yazalım:
\begin{equation}
    \rho(T) = \rho_0 + \left.\pd{\rho}{T}\right|_{T_0} (T - T_0) + \mathcal{O}\bigl((T-T_0)^2\bigr)
    \label{eq:rho_taylor}
\end{equation}

Termal genleşme katsayısı tanımı:
\begin{equation}
    \beta \equiv -\frac{1}{\rho_0}\left.\pd{\rho}{T}\right|_{T_0}
    \label{eq:beta_def}
\end{equation}

Eksi işareti konvansiyondur: çoğu akışkanda sıcaklık arttıkça yoğunluk \textit{azalır}, yani $\partial\rho/\partial T < 0$; böylece $\beta > 0$ olur.

Denklem~\eqref{eq:rho_taylor}'ye yerleştirelim:
\begin{equation}
    \rho(T) = \rho_0\bigl[1 - \beta(T - T_0)\bigr]
    \label{eq:boussinesq_eos}
\end{equation}

\begin{example}
Hava için $\beta \approx 1/T_0$ (ideal gaz). $T_0 = 300$~K'de $\beta = 1/300 \approx 3.33 \times 10^{-3}$~K$^{-1}$.
\begin{itemize}
    \item $\Delta T = 20$~K: $\beta\Delta T = 0.067 \ll 1$ $\checkmark$ (Boussinesq geçerli)
    \item $\Delta T = 50$~K: $\beta\Delta T = 0.167$ (sınırda)
    \item $\Delta T = 200$~K: $\beta\Delta T = 0.667$ (Boussinesq geçersiz!)
\end{itemize}
\end{example}

\textbf{Adım 2: Momentum denklemine yerleştirme.}

Tam (sıkıştırılabilir) momentum denklemi:
\begin{equation}
    \rho\DDt{\vect{u}} = -\gradient p + \mu\laplacian\vect{u} + \rho\vect{g} + \vect{F}
\end{equation}

Sol taraftaki $\rho$'yu Boussinesq'e göre $\rho_0$ ile değiştirin (çünkü $\Delta\rho/\rho_0 \ll 1$):
\begin{equation}
    \rho_0\DDt{\vect{u}} = -\gradient p + \mu\laplacian\vect{u} + \rho(T)\,\vect{g} + \vect{F}
\end{equation}

\textbf{Adım 3: Yerçekimi teriminin işlenmesi.}

$\rho(T)\vect{g}$ terimini açalım. $\vect{g} = -g\,\hat{\vect{e}}_y$ (aşağı yönde):
\begin{equation}
    \rho(T)\,\vect{g} = \rho_0[1 - \beta(T-T_0)](-g\,\hat{\vect{e}}_y) = -\rho_0 g\,\hat{\vect{e}}_y + \rho_0 g\beta(T-T_0)\hat{\vect{e}}_y
\end{equation}

İlk terim ($-\rho_0 g\hat{\vect{e}}_y$) hidrostatik basınçla dengelenir. \textit{Değiştirilmiş basınç} tanımlayalım:
\begin{equation}
    p' = p + \rho_0 g\,y
    \label{eq:modified_pressure}
\end{equation}

O zaman $-\gradient p - \rho_0 g\,\hat{\vect{e}}_y = -\gradient p'$ olur.

\textbf{Adım 4: Boussinesq momentum denklemi.}

\begin{equation}
    \rho_0\DDt{\vect{u}} = -\gradient p' + \mu\laplacian\vect{u} + \rho_0 g\beta(T - T_{\text{ref}})\hat{\vect{e}}_y + \vect{F}
\end{equation}

$\rho_0$'a bölersek:
\begin{equation}
    \boxed{\pd{\vect{u}}{t} + (\vect{u}\cdot\gradient)\vect{u} = -\frac{1}{\rho_0}\gradient p' + \nu\laplacian\vect{u} + g\beta(T - T_{\text{ref}})\hat{\vect{e}}_y + \frac{\vect{F}}{\rho_0}}
    \label{eq:boussinesq_momentum}
\end{equation}

\textbf{Kaldırma terimi} $g\beta(T - T_{\text{ref}})\hat{\vect{e}}_y$'nin yapısına dikkat edin:
\begin{itemize}
    \item $T > T_{\text{ref}}$: Akışkan sıcak $\to$ kuvvet yukarı ($+y$) $\to$ yükselir.
    \item $T < T_{\text{ref}}$: Akışkan soğuk $\to$ kuvvet aşağı ($-y$) $\to$ alçalır.
    \item $T = T_{\text{ref}}$: Kaldırma yok.
\end{itemize}

\textbf{Adım 5: Tam Boussinesq sistemi.}

Enerji denklemi değişmez (zaten $\rho = \rho_0$ ile türetildi):
\begin{subequations}
\label{eq:boussinesq_system}
\begin{align}
    \pd{\vect{u}}{t} + (\vect{u}\cdot\gradient)\vect{u} &= -\frac{1}{\rho_0}\gradient p' + \nu\laplacian\vect{u} + g\beta(T - T_{\text{ref}})\hat{\vect{e}}_y + \frac{\vect{F}}{\rho_0} \label{eq:bouss_mom} \\
    \pd{T}{t} + (\vect{u}\cdot\gradient)T &= \kappa\laplacian T \label{eq:bouss_ene} \\
    \divergence\vect{u} &= 0 \label{eq:bouss_cont}
\end{align}
\end{subequations}

\begin{remark}
Boussinesq'te $\divergence\vect{u} = 0$ \textbf{kesindir}. Bu denklem bir ``yaklaşım'' değildir: yoğunluk etkisi \textbf{yalnızca} $g\beta(T-T_{\text{ref}})$ teriminde yer alır. Süreklilik denklemi ``sanki sıkıştırılamaz'' gibidir. Bu, sayısal çözüm için büyük kolaylıktır çünkü basınç Poisson denklemiyle hesaplanabilir.
\end{remark}


% ============================================================
\section{Rayleigh--B\'{e}nard Konveksiyonu}
\label{sec:rayleigh_benard}

\subsection{Problem Tanımı}

Rayleigh--Bénard (RB) konveksiyonu, doğal konveksiyonun en temel ve en çok çalışılan modelidir:

\begin{itemize}
    \item Yatay iki plaka arasında akışkan.
    \item Alt plaka sıcak ($T_h$), üst plaka soğuk ($T_c$).
    \item Yerçekimi aşağı yönde.
    \item Dış zorlama yok ($\vect{F} = 0$, $U = 0$).
\end{itemize}

Tek kontrol parametresi $\Ray$ (ve $\Pran$).

\subsection{Lineer Stabilite Analizi}

Hidrostatik denge (akışkan hareketsiz) çözümü:
\begin{equation}
    \vect{u}_0 = 0, \qquad T_0(y) = T_h - \frac{\Delta T}{H}\,y, \qquad p_0(y) = p_{\text{atm}} + \rho_0 g (H-y)
    \label{eq:hydrostatic}
\end{equation}

Bu çözüm her zaman vardır. Soru şu: bu denge \textit{kararlı} mı?

Küçük bir pertürbasyon ekleyelim:
\begin{equation}
    \vect{u} = \vect{u}_0 + \vect{u}' = \vect{u}', \qquad T = T_0 + T', \qquad p = p_0 + p'
\end{equation}

Boussinesq denklemlerine yerleştirip $\mathcal{O}({\vect{u}'}^2)$ terimlerini ihmal edelim (lineerizasyon):
\begin{subequations}
\begin{align}
    \pd{\vect{u}'}{t} &= -\frac{1}{\rho_0}\gradient p' + \nu\laplacian\vect{u}' + g\beta\,T'\hat{\vect{e}}_y \label{eq:perturb_mom} \\
    \pd{T'}{t} - \frac{\Delta T}{H}v' &= \kappa\laplacian T' \label{eq:perturb_ene} \\
    \divergence\vect{u}' &= 0 \label{eq:perturb_cont}
\end{align}
\end{subequations}

Denklem~\eqref{eq:perturb_ene}'deki $-(\Delta T/H)v'$ terimi kilit terimdir: dikey hız pertürbasyonu $v'$ temel sıcaklık gradyanıyla etkileşerek sıcaklık pertürbasyonu üretir.

Normal mod analizi ($\vect{u}', T', p' \propto e^{\sigma t + i(k_x x + k_z z)}$) uygulanır ve $y$ yönünde uygun sınır koşulları ile bir \textbf{özdeğer problemi} elde edilir.

\subsection{Kritik Rayleigh Sayısı}

Lineer stabilite analizi sonucu:
\begin{equation}
    \boxed{\Ray_{\text{cr}} = \frac{27}{4}\pi^4 \approx 657.5 \quad \text{(free-slip)}}
    \label{eq:ra_cr_freeslip}
\end{equation}
\begin{equation}
    \boxed{\Ray_{\text{cr}} = 1707.76 \approx 1708 \quad \text{(no-slip)}}
    \label{eq:ra_cr_noslip}
\end{equation}

\begin{itemize}
    \item $\Ray < \Ray_{\text{cr}}$: Tüm pertürbasyonlar söner ($\sigma < 0$). Akışkan hareketsiz, $T$ doğrusal profilde kalır. Isı yalnızca iletimle taşınır: $\Nus = 1$.
    \item $\Ray = \Ray_{\text{cr}}$: Nötr kararlılık ($\sigma = 0$). Sonlu genlikli konveksiyon hücreleri oluşmaya başlar.
    \item $\Ray > \Ray_{\text{cr}}$: Pertürbasyonlar büyür ($\sigma > 0$). Konveksiyon başlar.
\end{itemize}

\subsection{Konveksiyon Yapıları}

$\Ray$ arttıkça akış yapıları zenginleşir:

\begin{enumerate}
    \item $\Ray \gtrsim \Ray_{\text{cr}}$: \textbf{Düzenli konveksiyon hücreleri.} 2B'de ``rolls'' (silindirler), 3B'de altıgen Bénard hücreleri. Sıcak akışkan yükselir, soğuyan akışkan alçalır.

    \item $\Ray \sim 10^4$--$10^5$: Hücreler kararsızlaşır. Zaman bağımlı davranış. Osilatörler, periyot-ikilemeleri.

    \item $\Ray \gtrsim 10^6$: \textbf{Türbülant konveksiyon.} Büyük ölçekli sirkulasyon (LSC) + türbülant plume'lar. Alt sınır tabakasından sıcak plume'lar kopup yükselir; üst sınır tabakasından soğuk plume'lar düşer.

    \item $\Ray \gtrsim 10^{10}$--$10^{14}$: ``Ultimate'' rejim tartışması. Sınır tabakaları kendileri türbülant mı oluyor? Açık araştırma sorusu.
\end{enumerate}


% ============================================================
\section{Zorlanmış Konveksiyon}
\label{sec:forced_convection}

Zorlanmış konveksiyonda akışkan hareketi dış kuvvetle sağlanır ve kaldırma kuvveti ihmal edilebilir ($\Rich \ll 1$).

Bu durumda sıcaklık denklemi \textbf{pasif skaler} gibi davranır:
\begin{subequations}
\begin{align}
    \pd{\vect{u}}{t} + (\vect{u}\cdot\gradient)\vect{u} &= -\frac{1}{\rho_0}\gradient p + \nu\laplacian\vect{u} + \frac{\vect{F}}{\rho_0} \\
    \pd{T}{t} + (\vect{u}\cdot\gradient)T &= \kappa\laplacian T \\
    \divergence\vect{u} &= 0
\end{align}
\end{subequations}

Dikkat: Momentum denkleminde $T$ yok. Hız alanı sıcaklıktan bağımsız çözülür; sonra sıcaklık dağılımı bilinen $\vect{u}$ ile bulunur. Bu tek yönlü bağlantıdır ve çözümü kolaylaştırır.

\subsection{Kolmogorov Forcing}

Bu projede kullanılan rüzgâr etkisi, \textbf{Kolmogorov forcing} ile modellenmiştir:
\begin{equation}
    \vect{F} = A\,\sin\!\left(k_f\,\frac{2\pi y}{L_y}\right)\hat{\vect{e}}_x
    \label{eq:kolmogorov_forcing}
\end{equation}

Burada:
\begin{itemize}
    \item $A$: Forcing genliği [m/s$^2$]
    \item $k_f$: Forcing dalga sayısı (genellikle $k_f = 1$)
    \item Kuvvet yalnızca $x$ yönünde ve $y$'ye sinüzoidal bağımlı
\end{itemize}

Neden bu forcing?
\begin{enumerate}
    \item \textbf{Analitik laminar çözüm vardır:} $u_{\text{lam}}(y) \propto \sin(k_f y)$. Bu, validasyon için kullanılabilir.
    \item \textbf{Enerji enjeksiyonu kontrollüdür:} Tek bir dalga sayısında enerji pompalanır, enerji kaskadı temiz çalışır.
    \item \textbf{Periyodik sınır koşullarıyla uyumludur:} $\sin$ fonksiyonu doğal olarak periyodiktir.
\end{enumerate}


% ============================================================
\section{Mixed Convection}
\label{sec:mixed_convection}

Mixed convection, zorlanmış ve doğal konveksiyonun \textit{birlikte} etkili olduğu rejimdir. $\Rich \sim \mathcal{O}(1)$ olduğunda ne kaldırma kuvveti ne de dış zorlama ihmal edilebilir.

\subsection{Fiziksel Yapılar}

\begin{figure}[H]
\centering
\includegraphics[width=\textwidth]{convection_types.png}
\caption{Üç konveksiyon rejimi: (sol) doğal konveksiyon ($\Rich \gg 1$, dikey plume'lar), (orta) zorlanmış konveksiyon ($\Rich \ll 1$, yatay akış), (sağ) mixed convection ($\Rich \sim 1$, eğik plume'lar ve karmaşık yapılar).}
\label{fig:convection_types}
\end{figure}

Mixed convection'da karmaşık akış yapıları ortaya çıkar:

\begin{enumerate}
    \item \textbf{Eğik plume'lar (tilted plumes):} Doğal konveksiyonda plume'lar dikey yükselir. Yatay rüzgâr onları yatırır. Plume açısı $\sim \arctan(\Rich^{1/2})$ ile ölçeklenir.

    \item \textbf{Kesilmiş konveksiyon hücreleri (sheared cells):} Bénard hücreleri rüzgâr tarafından deforme edilir. Rüzgâr yönünde uzar, dikey yönde sıkışır.

    \item \textbf{Büyük ölçekli sirkulasyon (LSC) + rüzgâr etkileşimi:} Doğal konveksiyonun yarattığı LSC ile dış akış etkileşir. Aynı yönde iseler güçlenir, ters yönde iseler zayıflar veya salınır.

    \item \textbf{Rejim geçişleri:} Parametre uzayında keskin olmayan geçişler. Küçük $\Delta\Ray$ değişimi akış yapısını tamamen değiştirebilir (histerezis mümkün).
\end{enumerate}

\subsection{Neden Zor?}

Mixed convection modellemesi zordur çünkü:
\begin{enumerate}
    \item \textbf{İki yönlü bağlantı:} Hız $\leftrightarrow$ sıcaklık. Bağımsız çözülemez.
    \item \textbf{Çoklu zaman ölçekleri:} Advektif zaman ölçeği $H/U$ ve kaldırma zaman ölçeği $H/\sqrt{g\beta\Delta T H}$ farklı olabilir.
    \item \textbf{Çoklu uzunluk ölçekleri:} Hız sınır tabakası, sıcaklık sınır tabakası, plume kalınlığı, forcing dalga boyu\ldots
    \item \textbf{3B etkileri güçlü:} 2B modeller yanıltıcı olabilir (spanwise instability, 3B plume etkileşimleri).
\end{enumerate}

\subsection{Literatürdeki Durum}

\begin{itemize}
    \item 2B RB konveksiyonu DNS ile kapsamlı çalışılmıştır ($\Ray$ = $10^6$--$10^{15}$).
    \item 3B RB konveksiyonu DNS ile $\Ray \lesssim 10^{10}$ kadar çalışılmıştır.
    \item Mixed convection (RB + forced) 2B'de sınırlı çalışma var.
    \item \textbf{3B mixed convection + neural operator: neredeyse çalışılmamış.} Bu, projemizin orijinal katkısıdır.
\end{itemize}


% ============================================================
\section{Non-Boussinesq Etkileri}
\label{sec:non_boussinesq}

\subsection{Boussinesq'in Sınırları}

Boussinesq yaklaşımı $\beta\Delta T \ll 1$ koşulunda geçerlidir. Bu koşul bozulduğunda:

\begin{enumerate}
    \item \textbf{Yoğunluk asimetrisi:} Sıcak ve soğuk plume'lar farklı davranır. Soğuk akışkan daha ağırdır, daha hızlı düşer. Sıcak akışkan daha hafiftir, daha yavaş yükselir. Bu asimetri Boussinesq'te yoktur.

    \item \textbf{Viskozite değişimi:} Gerçek akışkanlarda viskozite sıcaklıkla değişir. Hava için $\mu \propto T^{0.7}$. Bu da asimetriye katkıda bulunur.

    \item \textbf{Sıkıştırılamazlık bozulması:} $\divergence\vect{u} \neq 0$ olur. Stratifiye akışlarda bu önemli olabilir.
\end{enumerate}

\subsection{Ne Zaman Gerekir?}

\begin{center}
\begin{tabular}{rcl}
\toprule
$\beta\Delta T$ & \textbf{Durum} & \textbf{Yaklaşım} \\
\midrule
$< 0.1$ & Güvenli & Boussinesq yeterli \\
$0.1$--$0.2$ & Sınırda & Dikkatli değerlendirme gerekli \\
$> 0.2$ & Tehlikeli & Non-Boussinesq veya sıkıştırılabilir \\
\bottomrule
\end{tabular}
\end{center}

\begin{example}
Hava ($\beta = 1/T_0 = 1/300$~K$^{-1}$) için:
\begin{itemize}
    \item $\Delta T = 20$~K: $\beta\Delta T = 0.067$ $\to$ Boussinesq güvenle geçerli.
    \item $\Delta T = 50$~K: $\beta\Delta T = 0.167$ $\to$ Sınırda. Non-Boussinesq etkileri görülebilir.
    \item $\Delta T = 100$~K: $\beta\Delta T = 0.333$ $\to$ Boussinesq yetersiz.
\end{itemize}
Bu projede $\Delta T = 20$~K civarında çalışıyoruz, dolayısıyla Boussinesq geçerlidir.
\end{example}

\subsection{Non-Boussinesq Denklemler}

Tam non-Boussinesq formülasyon, yoğunluğu sıcaklığın fonksiyonu olarak tutar:
\begin{equation}
    \rho = \frac{\rho_0 T_0}{T} \quad \text{(ideal gaz)}
    \label{eq:non_bouss_eos}
\end{equation}

Bu durumda süreklilik denklemi:
\begin{equation}
    \pd{\rho}{t} + \divergence(\rho\vect{u}) = 0 \quad \text{ve genel olarak} \quad \divergence\vect{u} \neq 0
\end{equation}

Ancak düşük Mach sayısı ($\mathrm{Ma} \ll 1$) koşulunda ``low-Mach yaklaşımı'' kullanılabilir:
\begin{itemize}
    \item Akustik dalgalar filtrelenir (ses hızı $\to \infty$).
    \item Basınç, termodinamik ($p_0$) ve dinamik ($p'$) kısımlara ayrılır.
    \item $p_0$ uzayda sabit (kapalı kap) veya zamanda sabit (açık akış).
    \item Süreklilik denklemi: $\divergence\vect{u} = -(1/T)\,DT/Dt$ (sıfır değil!).
\end{itemize}

Non-Boussinesq modelleme, bu projenin ileri aşamasında (Boussinesq çalıştıktan sonra) bir genişletme hedefi olarak düşünülmektedir. INNATE mimarisine non-Boussinesq nöronları eklemek prensipte mümkündür.


% ============================================================
\section{Bölüm Özeti}
\label{sec:boussinesq_ozet}

\begin{enumerate}
    \item \textbf{Boussinesq yaklaşımı:} Yoğunluk değişimi yalnızca kaldırma teriminde. Geçerlilik: $\beta\Delta T \ll 1$.

    \item \textbf{Boussinesq denklemleri:}
    \begin{align*}
        \pd{\vect{u}}{t} + (\vect{u}\cdot\gradient)\vect{u} &= -\frac{1}{\rho_0}\gradient p + \nu\laplacian\vect{u} + g\beta(T-T_{\text{ref}})\hat{\vect{e}}_y + \frac{\vect{F}}{\rho_0} \\
        \pd{T}{t} + (\vect{u}\cdot\gradient)T &= \kappa\laplacian T \\
        \divergence\vect{u} &= 0
    \end{align*}

    \item \textbf{Rayleigh--Bénard:} Alttan ısıtma, üstten soğutma. $\Ray_{\text{cr}} \approx 1708$ (no-slip). $\Ray > \Ray_{\text{cr}}$ $\to$ konveksiyon.

    \item \textbf{Mixed convection:} Zorlanmış + doğal birlikte. $\Rich \sim \mathcal{O}(1)$. En karmaşık rejim.

    \item \textbf{Bu projede:} Boussinesq + Kolmogorov forcing + RB ısıtma $\to$ 3B mixed convection. $\Rey = 1000$--$10{,}000$, $\Ray = 10^5$--$10^7$, $\Pran = 0.71$.
\end{enumerate}
